\chapter{Conclusion}
\label{chap:conclusion}

\section{Summary of Findings}
\label{sec:concl-summary}
% Key results across all three evaluation levels
% Direct answer to the research question: how consistent are sovereign LLMs
%   in a structured therapeutic protocol -- and what does that mean clinically?

This thesis asked how consistent \gls{llm} therapeutic responses are
across independent runs within a structured \gls{irt} protocol. Three
evaluation methods were applied to ten models across 6{,}000 trials
spanning five temperature conditions.

The main findings are:

\begin{enumerate}
  \item \textbf{The \gls{eu}-sovereign flagship matches proprietary models
    at recommended temperatures.} \modelname{Mistral Large~3} achieves
    $J = 1.000$ at $T = 0.075$ and $T = 0.15$, equal to \modelname{GPT-5.4}
    and \modelname{Sonnet~4.6}. The sovereignty tradeoff in therapeutic
    stability is negligible when the model operates within its vendor-recommended
    range.

  \item \textbf{Instability is a decision-making problem, not an execution
    problem.} Temperature reduces plan consistency (Jaccard) but does not
    reduce plan-response alignment. Models select different strategies at
    higher temperatures but implement whichever strategy they select equally
    well.

  \item \textbf{Greedy decoding is not always optimal.} Three models achieve
    higher consistency at $T = 0.075$ than at $T = 0.0$. Whether this
    reflects a genuine benefit of low sampling variance or an artefact of
    the 20-trial sample size cannot be determined from these results alone.

  \item \textbf{The three evaluation levels capture distinct information.}
    Cross-method correlations are moderate at best ($\rho = 0.285$ to
    $0.560$). BERTScore is insensitive to plan-level changes. No single
    metric captures the full picture of therapeutic stability.

  \item \textbf{Patient profile matters.} A significant vignette effect
    ($p = 0.007$) means that stability evaluations based on aggregate
    metrics alone may miss per-profile weaknesses.
\end{enumerate}

Regarding the two contributions stated in \autoref{sec:intro-objectives}:
the three-level framework (Contribution~1) fulfilled its design goal by
identifying that instability is concentrated at the strategy selection
layer, a finding that neither BERTScore nor alignment alone could have
revealed. The cross-model comparison (Contribution~2) demonstrates that
the \gls{eu}-sovereign \modelname{Mistral Large~3} achieves stability
comparable to proprietary models at its recommended operating temperature,
supporting its viability for clinical deployment under sovereignty
constraints.


\section{Limitations}
\label{sec:concl-limitations}
% Scope: single protocol, single phase (IRT rescripting only)
% Simulation: synthetic patients differ from real clinical interaction in
%   ways that cannot be fully controlled
% Six vignettes sample the difficulty range but cannot claim exhaustive
%   coverage of patient variability
% Method 3 is exploratory -- judge reliability is unverified and should
%   not be treated as ground truth
% Fixed taxonomy constrains the measurable strategy space -- unmeasured
%   strategies may show different stability patterns
% Computational stability does not equal clinical safety

The study evaluates a single phase (rescripting) of a single protocol
(\gls{irt}). The rescripting phase was selected because it requires
active strategy decisions, but stability in other stages or other
therapeutic protocols may differ.

All patient interactions are simulated. Synthetic patients follow scripted
behavioural profiles and do not exhibit the unpredictable responses of
real patients. Six profiles sample the difficulty range
(\autoref{sec:meth-dialogue}) but cannot claim exhaustive coverage. No
human participants were involved in this study, and the results measure
computational consistency rather than clinical safety or efficacy.
\gls{irt} targets nightmares rather than psychiatric diagnoses, which
limits clinical risk, but any deployment to real patients would require
validation beyond what this framework provides.

The \gls{llm} judge (\modelname{Gemini 3.1 Pro}) has only been validated
on a small subset of trials due to the volume of data (6{,}000 trials).
Method~3 results should be interpreted as a computational estimate rather
than a clinically validated measure.

The six-category strategy taxonomy constrains what can be measured.
Strategies that fall outside these categories are invisible to Method~1,
and different category boundaries would produce different absolute
consistency scores.

Finally, consistency is not the same as quality. A model that consistently
selects an inappropriate strategy scores high on plan consistency.
Stability is necessary for clinical deployment but not sufficient on its
own.


\section{Future Work}
\label{sec:concl-future}
% Full IRT protocol evaluation -- beyond the rescripting phase
% Extension to other manualized therapies (CBT, DBT, PE)
% Longitudinal drift: does stability degrade across model version updates?
% Multi-language evaluation -- German therapeutic context for reDreamAI
% Clinical validation: human therapist ratings as external ground truth
% Replication with future model generations to track whether stability
%   improves across release cycles

Several directions follow from this work.

Extending the evaluation beyond the rescripting phase to all five
\gls{irt} stages would test whether the observed patterns hold when the
model faces different types of clinical decisions. More broadly, the
three-level framework applies to other therapy protocols where the
therapist chooses from a defined set of techniques, such as \gls{cbt} or
motivational interviewing. Each would need its own strategy taxonomy, but
the evaluation setup carries over.

Since model vendors update weights and infrastructure without notice,
repeating the evaluation across model versions would show whether
stability improves or fluctuates over time. reDreamAI targets
German-speaking users, so evaluating stability in German would test
whether the English-language patterns transfer.

On the measurement side, a larger human annotation study against the same
ternary rubric would strengthen the judge validation that this study could
only perform on a small subset. Comparing human and \gls{llm} judge
scores would clarify how much trust to place in automated alignment
assessment.

Finally, newer model architectures may affect output consistency
independently of temperature. \modelname{Qwen~3.5} already uses Gated
DeltaNet \parencite{YJTEGJH4}, a hybrid attention mechanism where gated
memory updates control how much prior context is retained or forgotten.
In principle, this could reduce the effect where an early token divergence
propagates through the rest of the response. The results here do not show
a clear stability advantage for \modelname{Qwen~3.5} over standard
transformers, but whether architectural choices can improve therapeutic
consistency beyond what temperature tuning achieves remains open.
